\rhead{MN-CIV-A: AI \& Big Data for Urban Systems}
\phantomsection
\section*{AI \& Big Data for Urban Systems (MN-CIV-A)}
\label{minor:CIV_L2}

\noindent \small{\hyperref[main:scheme]{$\leftarrow$ Back to Scheme}} \vspace{0.5cm}

\noindent \textbf{Credits:} 2 (2-0-0)

\subsection*{Course Content}
\noindent\textbf{Unit 1} \\
\textit{Urban Big Data Engineering and Platforms:} The velocity-volume-variety challenge in urban data: millions of connected devices generating continuous streams across heterogeneous modalities at city scale; Lambda and Kappa architectures for urban data pipelines: batch layer (Hadoop, Spark) for historical reprocessing vs.\ speed layer (Kafka, Flink) for real-time analytics; Urban data lakes: organizing raw, processed, and curated data tiers in object storage with schema-on-read; Apache Spark for distributed urban analytics: RDD and DataFrame abstractions, spatial UDFs with Sedona, and partitioning strategies for geospatial workloads; Stream processing for urban event detection: sliding and tumbling window aggregations on traffic, utility, and environmental sensor streams; Data quality at scale: automated schema validation, outlier flagging, and lineage tracking in urban data pipelines as a software engineering discipline; Federated urban data ecosystems: data spaces, sovereign data sharing (Gaia-X, FIWARE), and the IDSA connector architecture enabling cross-city data exchange without centralization.

\vspace{0.3cm}

\noindent\textbf{Unit 2} \\
\textit{Spatio-Temporal Prediction and Urban Forecasting:} Urban forecasting as a spatio-temporal regression problem: demand, flow, and environmental variables as functions of both location and time; Spatio-temporal kriging and covariance separability as the geostatistical baseline for urban prediction; Recurrent architectures for urban time series: LSTM and GRU for multivariate sensor forecasting with spatial covariates; Graph Convolutional Networks (GCN) for network-structured urban data: diffusion convolution on the road graph adjacency matrix for traffic speed and flow forecasting; Spatial transformer networks for demand forecasting: treating the city grid as an image and learning spatial dependencies via convolutional filters; Uncertainty quantification in urban forecasts: conformal prediction intervals for traffic ETAs and energy demand as inputs to risk-aware operational decisions; Hierarchical forecasting: reconciling predictions at neighborhood, district, and city scales using bottom-up and optimal reconciliation methods (MinT).

\vspace{0.3cm}

\noindent\textbf{Unit 3} \\
\textit{Reinforcement Learning for Urban Control:} Urban control problems as Markov Decision Processes: the state space, action space, reward function, and transition dynamics for traffic signal control, energy dispatch, and shared mobility rebalancing; Multi-agent reinforcement learning (MARL) for city-scale control: decentralized execution with centralized training (CTDE) for coordinating hundreds of intersections simultaneously; Deep Q-Network (DQN) and Proximal Policy Optimization (PPO) applied to adaptive traffic signal control: outperforming Webster's fixed-time plans in simulation benchmarks (CityFlow, SUMO); Demand-responsive transit: RL-based dynamic routing and stop insertion for on-demand bus services; Smart grid demand response as a multi-period stochastic optimization: model predictive control vs.\ RL for real-time load scheduling under renewable uncertainty; Sim-to-real transfer challenges in urban RL: the distribution shift between calibrated city simulators and the physical complexity of real urban environments.

\vspace{0.3cm}

\noindent\textbf{Unit 4} \\
\textit{Urban AI Applications: Vision, NLP, and Multimodal Sensing:} Computer vision for urban intelligence: vehicle and pedestrian detection from CCTV using YOLO-family detectors; crowd density estimation via density map regression; anomaly detection for public safety; Street-level imagery analytics: semantic segmentation of Mapillary and Google Street View imagery for automated urban audit (sidewalk quality, green cover, facade condition) as a scalable alternative to field surveys; NLP for urban governance: topic modeling (LDA, BERTopic) on 311 service requests and social media to extract citizen-reported infrastructure issues and sentiment toward city services; Multimodal fusion for urban situational awareness: combining satellite imagery, street-level video, social media text, and sensor telemetry in a unified representation for emergency response coordination; Noise and air quality mapping: fusing sparse sensor networks with land-use regression models and deep learning to produce high-resolution pollution maps at city scale.

\vspace{0.3cm}

\noindent\textbf{Unit 5} \\
\textit{Urban Simulation, Digital Twins, and the Future City:} Agent-based models (ABMs) as the simulation paradigm for emergent urban phenomena: agents, behaviors, environments, and the MATSim transport simulation as a city-scale activity-based mobility model; Urban digital twins: the city-scale extension of asset-level twins integrating real-time sensor data, 3D city models (CityGML, 3D Tiles), and simulation engines for scenario planning and operational monitoring; Generative AI for urban design: diffusion models and GANs conditioned on zoning constraints and demographic targets for synthesizing urban layouts, building massing, and land use configurations; Equity-aware urban AI: algorithmic fairness metrics (demographic parity, equalized odds) applied to urban resource allocation, and participatory AI methods for incorporating community input into model design; Autonomous urban systems: self-driving vehicle fleet management, drone delivery network optimization, and robotic last-mile logistics as the next wave of AI integration into urban infrastructure; Climate-adaptive cities: using big data and AI for urban heat island mitigation, flood early-warning systems, and optimizing green infrastructure placement as responses to climate-driven urban risk.

\newpage
\rhead{Curriculum Schema}